\documentclass[11pt]{article}
\usepackage[margin=0.85in]{geometry}
\usepackage{amsmath,hyperref}
\title{Separating wetland emissions and OH with annual GOSAT inversions}
\author{Reading note on Qu et al. (2024)}\date{}
\begin{document}\maketitle
\noindent\textit{Inverse modeling of 2010--2022 satellite observations
shows that inundation of the wet tropics drove the 2020--2022 methane surge.}
PNAS 121, e2402730121.
\href{https://doi.org/10.1073/pnas.2402730121}{doi:10.1073/pnas.2402730121}.
The main article and supplementary Figs.~S1--S5 and Table~S1 were obtained.

\paragraph{Method.}
Independent annual analytical inversions use GOSAT column methane and
GEOS-Chem 12.5.0, driven by MERRA2 at $4^\circ\times5^\circ$.
Each inversion has 1009 annual nonwetland grid-cell fluxes, 168 monthly
wetland fluxes over 14 regions, and two annual hemispheric OH factors.
Land and ocean-glint retrievals are retained between $60^\circ$S and
$60^\circ$N. Every year uses the same prior emissions and a global
initial-concentration scaling based on the first ten days of GOSAT
observations. The plotted continuous model concentration trajectory
does not apply those yearly resets. Priors use GFEIv2, EDGARv4.3.2,
and nine selected WetCHARTs members. A ten-member ensemble varies
prior uncertainty and regularization one parameter at a time.

\paragraph{Findings.}
Table~S1 reports emissions of 550, 587, 590, and 569 Tg yr$^{-1}$
in 2019--2022: anomalies of 0, 37, 40, and 19 relative to 2019.
OH removal is 478, 467, 468, and 453 Tg yr$^{-1}$; other removal
is 75, 75, 74, and 79 Tg yr$^{-1}$. These yield model growth of
$-3$, 45, 48, and 37 Tg yr$^{-1}$, respectively. In particular,
the negative modeled 2019 growth cannot be silently replaced with
NOAA's positive 2019 growth. These published numbers are rounded to
integer Tg yr$^{-1}$. Mean OH lifetime over 2010--2022 is 10.9 years.
The paper attributes 28\% of the surge to OH and emphasizes emission
increases in Equatorial Asia and Africa (43\% and 30\% of the
emission surge). Its earlier South American increase contributes little
specifically to 2020--2022.

\paragraph{Why the comparison is informative.}
This inversion separately optimizes monthly regional wetlands, whereas
He et al. optimizes annual grid-cell net flux and later partitions it
using prior sector fractions. It also resets to a fixed prior each year
instead of propagating the previous posterior. These differences alter
both sector attribution and the 2019 reference state. Tropical
water-storage correlations support an inundation interpretation but do
not uniquely establish causation; GRACE measures relative storage,
not inundated wetland area. The paper itself notes discrepancies with
livestock inventories and uncertainty in precipitation responses.

\paragraph{Data and adjustment decision.}
Table~S1 supports an exact transcription of rounded global fluxes.
Supplementary Fig.~S3 is area-weighted OH, which must not be substituted
for methane-reaction-weighted loss frequency. The smaller OH removal
fluxes must first be converted using matched methane burdens before
applying He's $-m_0\Delta k_{\mathrm{OH}}$ term. Other sinks combine
soil and chemical loss, so splitting them requires additional data.
We retain those limitations instead of inventing missing components.
\end{document}
